\section{Running the verifications}\label{sec:verification}

The data and code for the two computer-assisted steps, \cref{lem:cancellation}
and the Sturm-bound computation in the proof of \cref{thm:K5-recurrence}, are
described in \cref{app:certificate,app:verifier,app:sturm} and are available in
the repository
\begin{center}
\url{https://github.com/ramanujamath/non-rascoe-partitions-mod4}
\end{center}
in the directory \texttt{verification/appendix}. Download the three programs into one directory:
\begin{itemize}
\item \texttt{decode\_certificate.py},
\item \texttt{verify\_original\_certificate.py},
\item \texttt{prove\_square\_scaling\_29.py},
\end{itemize}
and run them, with Python $3.10$ or later and without the optimization flag
\texttt{-O}:
\begin{verbatim}
python3 decode_certificate.py
python3 verify_original_certificate.py
python3 prove_square_scaling_29.py
\end{verbatim}
Only the Python standard library is needed. The first program reconstructs the
certificate and checks its SHA-256 digest; the second reconstructs the
signatures of \cref{lem:cancellation} and checks every required identity (about
seven seconds); the third computes $\mathfrak F$ modulo $2$ through the Sturm
bound (about one second). Neither verification uses the search procedure that
produced the certificate, or any external data.
